%%%%%%%%%%%%%%%%%%%%%%%%%%%%%%%%%%%%%%%%%%%%%%%%%%%%%%%%%%%%%%%%%%%%%%%%%%%
%
% Generic template for TFC/TFM/TFG/Tesis
%
% $Id: estudioTeorico.tex,v 1.5 2015/06/05 00:05:19 macias Exp $
%
% By:
%  + Javier Macías-Guarasa.
%    Departamento de Electrónica
%    Universidad de Alcalá
%  + Roberto Barra-Chicote.
%    Departamento de Ingeniería Electrónica
%    Universidad Politécnica de Madrid
%
% Based on original sources by Roberto Barra, Manuel Ocaña, Jesús Nuevo, Pedro Revenga, Fernando Herránz and Noelia Hernández. Thanks a lot to all of them, and to the many anonymous contributors found (thanks to google) that provided help in setting all this up.
%
% See also the additionalContributors.txt file to check the name of additional contributors to this work.
%
% If you think you can add pieces of relevant/useful examples, improvements, please contact us at (macias@depeca.uah.es)
%
% You can freely use this template and please contribute with comments or suggestions!!!
%
%%%%%%%%%%%%%%%%%%%%%%%%%%%%%%%%%%%%%%%%%%%%%%%%%%%%%%%%%%%%%%%%%%%%%%%%%%%

\chapter{Elementos básicos de \LaTeX{}}
\label{cha:elementos-basicos}

\begin{FraseCelebre}
  \begin{Frase}
    Y así, del mucho leer y del poco dormir, se le secó el cerebro de manera que vino a perder el juicio\footnote{Tomado de ejemplos del proyecto \texis{}.}.
  \end{Frase}
  \begin{Fuente}
    Miguel de Cervantes Saavedra
  \end{Fuente}
\end{FraseCelebre}


\section{Estructura, listas y citas}
\label{sec:estructura-listas-citas}

Este capítulo reúne ejemplos breves de elementos habituales de \LaTeX{}. Puedes consultar directamente sus fuentes y adaptar únicamente los fragmentos que necesites.

Una lista sin numerar puede construirse mediante el entorno \texttt{itemize}:

\begin{itemize}
  \item Primer elemento de ejemplo.
  \item Segundo elemento de ejemplo.
\end{itemize}

Las fuentes bibliográficas se citan con \texttt{\textbackslash{}cite}; por ejemplo, esta plantilla describe el proyecto \texis{} en \cite{texis}. También puedes añadir una aclaración mediante una nota al pie\footnote{Este es un ejemplo de nota al pie.}.

\section{Matemáticas y referencias cruzadas}
\label{sec:matematicas-referencias}

Puedes incluir expresiones matemáticas en línea, como $\sigma=0.75$, o utilizar un entorno independiente para una ecuación numerada:

\begin{equation}
  \label{eq:ejemplo-basico}
  p[q_t=\sigma_t\mid q_{t-1}=\sigma_{t-1}]
\end{equation}

La ecuación~\ref{eq:ejemplo-basico} muestra cómo asociar una etiqueta mediante \texttt{\textbackslash{}label} y recuperarla con \texttt{\textbackslash{}ref}. El mismo mecanismo se utiliza para capítulos, secciones, figuras, tablas, algoritmos y fragmentos de código.

\section{Niveles de organización}
\label{sec:niveles-organizacion}

A continuación se muestran los niveles de división disponibles por debajo de un capítulo. Su aparición en el índice depende del valor configurado para \texttt{tocdepth}.


\subsection{Ejemplo de subsección}
\label{sec:subseccion}


\subsubsection{Ejemplo de subsubsección}
\label{sec:subsubseccion}

\paragraph{Ejemplo de párrafo}
\label{sec:paragraph-1}


\subparagraph{Ejemplo de subpárrafo}
\label{sec:subparagraph}




%%%%%%%%%%%%%%%%%%%%%%%%%%%%%%%%%%%%%%%%%%%%%%%%%%%%%%%%%%%%%%%%%%%%%%%%%%%
%
% Generic template for TFC/TFM/TFG/Tesis
%
% $Id: resultados.tex,v 1.7 2016/03/31 10:44:23 macias Exp $
%
% By:
%  + Javier Macías-Guarasa.
%    Departamento de Electrónica
%    Universidad de Alcalá
%  + Roberto Barra-Chicote.
%    Departamento de Ingeniería Electrónica
%    Universidad Politécnica de Madrid
%
% Based on original sources by Roberto Barra, Manuel Ocaña, Jesús Nuevo, Pedro Revenga, Fernando Herránz and Noelia Hernández. Thanks a lot to all of them, and to the many anonymous contributors found (thanks to google) that provided help in setting all this up.
%
% See also the additionalContributors.txt file to check the name of additional contributors to this work.
%
% If you think you can add pieces of relevant/useful examples, improvements, please contact us at (macias@depeca.uah.es)
%
% You can freely use this template and please contribute with comments or suggestions!!!
%
%%%%%%%%%%%%%%%%%%%%%%%%%%%%%%%%%%%%%%%%%%%%%%%%%%%%%%%%%%%%%%%%%%%%%%%%%%%

\section{Figuras y tablas}
\label{sec:figuras-tablas}


\subsection{Figuras de uso habitual}
\label{sec:figuras-habituales}

Este capítulo presenta ejemplos de figuras y tablas que pueden adaptarse a la mayoría de los documentos. Los casos que requieren una composición más elaborada se reúnen en el capítulo~\ref{cha:ejemplos-avanzados}.

La figura~\ref{fig:figura-convencional} muestra la inclusión habitual de una imagen con \texttt{figure}, \texttt{\textbackslash{}includegraphics}, un pie descriptivo y una etiqueta para citarla desde el texto.

\begin{figure}[h]
  \centering
  \includegraphics[width=0.55\textwidth]{Figure1}
  \caption{Ejemplo de figura convencional.}
  \label{fig:figura-convencional}
\end{figure}

\subsection{Tabla convencional}
\label{sec:tabla-convencional}

La tabla~\ref{tab:table1} muestra una estructura convencional con título, etiqueta y referencia cruzada.

\begin{table}
  % increase table row spacing, adjust to taste
  \renewcommand{\arraystretch}{1.3}
  \caption{Comparativa.}
  \label{tab:table1}
  \begin{center}
    % Some packages, such as MDW tools, offer better commands for making tables than the plain LaTeX2e tabular which is used here.
    \begin{tabular}{|c|c|c|}
      \hline
      Method & Training Time & Man-Work (\%)\\
      \hline
      Propagation model & $<$ 30 sec & 5\\
      \hline
      Manual & 9 h 30 min & 24\\
      \hline
      Automatic & 2 h & 10 8\\
      \hline
    \end{tabular}
  \end{center}
\end{table}

\subsection{Subfiguras}
\label{sec:subfiguras}

El entorno \texttt{subfigure} permite reunir varias imágenes dentro de una misma figura. La figura~\ref{fig:fig3} contiene dos subfiguras, que también pueden citarse individualmente; por ejemplo, \ref{fig:fig3}.\subref{fig:fig3b}.

% For this to work you need to (in preamble.tex):
% - remove \usepackage{subfig}
% - add \usepackage{caption}
% - add \usepackage{subcaption}
\begin{figure}
  \centering
  \begin{subfigure}[b]{0.3\textwidth}
    \includegraphics[width=\textwidth]{Figure2}
    \caption{Mean Entropy.}
    \label{fig:fig3a}
  \end{subfigure}%
  ~ %add desired spacing between images, e. g. ~, \quad, \qquad etc.
  % (or a blank line to force the subfigure onto a new line)
  \begin{subfigure}[b]{0.3\textwidth}
    \includegraphics[width=\textwidth]{Figure3}
    \caption{Error Percentage.}
    \label{fig:fig3b}
  \end{subfigure}
  \caption{Optimal trames number in the training data set.}
  \label{fig:fig3}
\end{figure}

Los ejemplos avanzados de figuras y tablas, incluidas las tablas multipágina y apaisadas, las figuras rodeadas por texto y las composiciones con numerosas subfiguras, se encuentran en el capítulo~\ref{cha:ejemplos-avanzados}.


%%% Local Variables:
%%% TeX-master: "../book"
%%% End:


%%% Local Variables:
%%% TeX-master: "../book"
%%% End:
